\chapter[Metodologia]{Metodologia}

Este capítulo descreve a metodologia adotada para o desenvolvimento deste trabalho, descrevendo as etapas realizadas desde a obtenção dos dados 
até a avaliação do desempenho dos modelos propostos. Inicialmente, foi realizada a seleção e organização da MIT-BIH Arrhythmia 
Database, seguida pelas etapas de pré-processamento dos sinais, depois os dados foram estruturados em conjuntos de treinamento, validação e 
teste. Na sequência, foram projetadas e implementadas duas arquiteturas de redes neurais convolucionais: uma CNN1D, aplicada diretamente aos 
sinais temporais de ECG, e uma CNN2D, na qual os batimentos foram representados em formato bidimensional. Ambas as arquiteturas foram treinadas sob a 
mesma base de dados, sendo os parâmetros de configuração e arquitetura ajustados de forma empírica por meio de experimentações. Essa padronização 
metodológica possibilitou uma avaliação e comparação direta de desempenho.

\section{Base de dados MIT-BIH}

A base de dados utilizada neste trabalho foi a MIT-BIH Arrhythmia Database, um dos conjuntos de dados mais empregados em 
pesquisas relacionadas à detecção e classificação de arritmias cardíacas. Essa base foi desenvolvida pelo MIT em parceria com o Beth Israel 
Hospital, com o objetivo de fornecer um padrão de referência para avaliação de algoritmos de análise 
automática de sinais de ECG.

A base do MIT-BIH é composta por 48 registros de ECG, com aproximadamente 30 minutos cada, obtidos a partir de pacientes 
monitorados em ambiente clínico. Os sinais foram adquiridos com frequência de amostragem de 360 Hz. Cada registro foi anotado 
por especialistas, indicando a localização temporal dos batimentos e sua respectiva classificação.

Existem 15 tipos de classes no banco que podem ser vista na Tabela 1, como esperado em qualquer base de dados médicos a quantidade de casos 
normais é muito supeior as demais.

\begin{table}[ht]
  \caption{Classes da base de dados do MIT-BIH}
  \label{tab: table1}
  \centering

  \begin{tabular}{cc}
    \hline
    \bfseries Classe & \bfseries Símbolo\\
    \hline
    Normal beat & N\\
    Left bundle branch block beat & L\\
    Right bundle branch block beat & R\\
    Atrial premature beat & A\\
    Aberrated atrial premature beat & a\\
    Nodal (junctional) premature beat & J\\
    Supraventricular premature beat & S\\
    Premature ventricular contraction & V\\
    Fusion of ventricular and normal beat & F\\
    Atrial escape beat & e\\
    Nodal (junctional) escape beat & j\\
    Ventricular escape beat & E\\
    Fusion of paced and normal beat & f\\
    Paced beat & /\\
    Unclassifiable beat & Q\\
    \hline
  \end{tabular}

  \source
\end{table}

Para reduzir a grande diferença entre as classes foi utilizado o padrão da Association for the Advancement of Medical Instrumentation (AAMI), 
uma organização que é líder em padrões para tecnologia de saúde. O padrão da AAMI visa agrupar as 15 classes da base do MIT-BIH em 5 grupos 
com características parecidas, como pode ser visto na Tabela 2

\begin{table}[ht]
  \caption{Classificação das classes de ECG segundo o padrão AAMI}
  \label{tab: table2}
  \centering
  \begin{tabular}{cc}
    \hline
    \bfseries Classe AAMI & \bfseries Tipos de batimentos (MIT-BIH) \\
    \hline
    Normal (N) & N, L, R, e, j \\
    Supraventricular (S) & A, a, J, S \\
    Ventricular (V) & V, E \\
    Fusion (F) & F \\
    Unknown (Q) & /, f, Q \\
    \hline
  \end{tabular}

  \legend{Adaptado de \cite{DasMK}}
\end{table}


\section{Pré-processamento}

A etapa de pré-processamento teve como objetivo preparar os sinais do ECG para utilização no treinamento e avaliação dos 
modelos de redes neurais convolucionais. Essa etapa é fundamental, uma vez que os sinais brutos apresentam variações de duração, amplitude e 
distribuição entre as classes, o que pode comprometer o desempenho.

Inicialmente, foram selecionados todos os registros da base MIT-BIH Arrhythmia Database, compostos por sinais de ECG adquiridos com frequência de 
amostragem de 360 Hz. Para cada registro, foi considerada apenas uma derivação do sinal, correspondente ao canal 0, que representa a derivação MLII. 
As anotações disponibilizadas pela base foram utilizadas para identificar a posição temporal dos batimentos cardíacos e seus respectivos rótulos.

A partir das anotações, cada batimento foi segmentado em uma janela centrada no pico do complexo QRS. Para essa segmentação, foi definida uma 
janela inicial maior, com 400 amostras, de modo a permitir deslocamentos temporais durante o processo de aumento de dados. Batimentos localizados 
próximos às extremidades do sinal, que não permitiam a extração completa da janela, foram descartados.

Em seguida, cada segmento foi normalizado individualmente por meio da normalização min–max, garantindo que os valores do sinal estivessem 
restritos ao intervalo entre 0 e 1. Essa normalização reduz a influência de variações de amplitude entre diferentes registros e 
contribui para a estabilidade do treinamento das redes neurais.

Após a segmentação e normalização, os batimentos foram mapeados para as cinco classes definidas pelo padrão AAMI: batimentos normais (N), 
supraventriculares (S), ventriculares (V), de fusão (F) e desconhecidos (Q). Esse mapeamento foi realizado a partir das classes originais da 
base MIT-BIH, conforme recomendado na literatura, possibilitando a comparação dos resultados com outros trabalhos da área.

Devido ao desbalanceamento natural entre as classes da base MIT-BIH Arrhythmia Database, foi aplicado um processo de balanceamento no conjunto de 
treinamento. A definição das quantidades por classe foi baseada na classe minoritária, que possui aproximadamente 800 amostras no total. Dessa forma, 
foram reservadas 200 amostras para validação e 200 para teste, restando cerca de 400 amostras para o treinamento nessa classe. Para garantir equilíbrio 
entre as classes durante o treinamento, foi estabelecido um total fixo de 1600 amostras por classe.

Classes com quantidade superior a esse valor foram submetidas à subamostragem aleatória. Já as classes com menor quantidade de amostras foram balanceadas 
por meio de aumento artificial de dados até atingir 1600 amostras por classe. Esse processo foi realizado a partir da replicação aleatória de amostras 
existentes, seguida da aplicação de transformações estocásticas.

Foram empregadas três técnicas de aumento de dados. A primeira consiste no deslocamento temporal da janela de segmentação, no qual um segmento expandido 
de 400 amostras (360 amostras centrais mais uma margem de ±20 amostras) é recortado aleatoriamente, gerando variações na posição do batimento dentro da 
janela final. A segunda técnica corresponde à adição de ruído gaussiano com média zero e desvio padrão de 0,01, aplicado ao sinal normalizado. A terceira 
técnica consiste na variação de escala da amplitude do sinal, por meio de um fator multiplicativo aleatório uniformemente distribuído no intervalo de 
0,95 a 1,05.

Para cada nova amostra gerada, uma das três técnicas de aumento de dados foi selecionada aleatoriamente com igual probabilidade. O número de amostras 
geradas por técnica não foi previamente fixado, sendo determinado de forma dinâmica de acordo com a necessidade de balanceamento de cada classe, até 
atingir o total de 1600 amostras no conjunto de treinamento. Ao final desse processo, os conjuntos gerados foram armazenados em arquivos, a fim de 
evitar a necessidade de reprocessamento em execuções futuras.

Após o balanceamento, todos os segmentos foram recortados para uma janela final de 360 amostras, correspondente a aproximadamente um segundo de sinal, 
centrada no batimento cardíaco. Para os conjuntos de validação e teste, não foi aplicado aumento de dados, sendo realizada apenas a extração central da 
janela, de modo a preservar a fidelidade dos sinais utilizados na avaliação.

\begin{figure}[ht]
  \caption{Exemplo de batimento da classe N}
  \label{fig: fig15}
  \centering
  \includegraphics[width=0.7\textwidth]{Exemplo_Classe_N.png}
  \source
\end{figure}

Por fim, os conjuntos finais de treinamento e teste foram armazenados em arquivos no formato .npy, permitindo reutilização eficiente dos dados
nas etapas subsequentes de treinamento e avaliação dos modelos.

\section{Arquitetura e treinamento da CNN1D}

Nesta etapa foi desenvolvido e treinado um modelo de CNN1D, aplicado diretamente aos sinais temporais de eletrocardiograma previamente 
segmentados e normalizados. A utilização de CNN1D é adequada para esse tipo de dado, pois permite a extração automática de características 
locais ao longo do eixo temporal do sinal, como padrões relacionados ao complexo QRS, ondas P e T.

A arquitetura proposta recebeu como entrada segmentos de ECG com 360 amostras, representando um único canal temporal. O modelo foi estruturado 
de forma sequencial, composto por três blocos convolucionais seguidos por uma etapa de classificação totalmente conectada.

O primeiro bloco convolucional é formado por uma camada Conv1D com 32 filtros e tamanho de kernel igual a 5, utilizando a função de ativação 
ReLU. Em seguida, foi aplicada a normalização em lote (Batch Normalization), com o objetivo de estabilizar o processo de treinamento e acelerar
a convergência do modelo. Posteriormente, uma camada de MaxPooling1D com fator de redução igual a 2 foi utilizada para diminuir a 
dimensionalidade do sinal e reduzir a complexidade computacional.

O segundo bloco segue estrutura semelhante, contendo uma camada Conv1D com 64 filtros e kernel de tamanho 5, seguida novamente por Batch 
Normalization e MaxPooling1D. Esse bloco permite a extração de padrões temporais mais complexos e hierárquicos presentes nos sinais de ECG.

O terceiro bloco convolucional é composto por uma camada Conv1D com 128 filtros e kernel de tamanho 3, também utilizando ativação ReLU. Assim 
como nos blocos anteriores, foram aplicadas as etapas de normalização em lote e pooling, aprofundando a capacidade de representação do 
modelo e reduzindo progressivamente a dimensão do sinal.

Após os blocos convolucionais, a saída do modelo foi transformada em um vetor unidimensional por meio da camada Flatten. Em seguida, foi 
aplicada uma camada de Dropout com taxa de 50\%, com o objetivo de reduzir o risco de sobreajuste durante o treinamento. A etapa de classificação
é composta por uma camada densa com 100 neurônios e função de ativação ReLU, seguida por uma camada de saída com cinco neurônios e função de 
ativação Softmax, correspondente às cinco classes definidas pelo padrão AAMI.

\begin{table}[ht]
  \caption{Detalhamento da arquitetura proposta}
  \label{tab: table3}
  \centering

  \begin{tabular}{cccc}
    \hline
    \bfseries Camada & \bfseries Filtros/Neurônios & \bfseries Tamanho do Kernel & \bfseries Saída\\
    \hline
    Entrada	&	-	&	-	&	(360, 1)\\
    Conv1D 1	&	32	&	5	&	(356, 32)\\
    Batch Norm 1	&	-	&	-	&	(356, 32)\\
    Max Pooling 1	&	-	&	2	&	(178, 32)\\
    Conv1D 2	&	64	&	5	&	(174, 64)\\
    Batch Norm 2	&	-	&	-	&	(174, 64)\\
    Max Pooling 2	&	-	&	2	&	(87, 64)\\
    Conv1D 3	&	128	&	3	&	(85, 128)\\
    Batch Norm 3	&	-	&	-	&	(85, 128)\\
    Max Pooling 3	&	-	&	2	&	(42, 128)\\
    Flatten	&	-	&	-	&	(5376)\\
    Dropout	&	-	&	-	&	(5376)\\
    Densa (FC)	&	100	&	-	&	(100)\\
    Saída (Softmax)	&	5	&	-	&	(5)\\

    \hline
  \end{tabular}

  \source
\end{table}

O modelo foi treinado utilizando o otimizador Adam, com taxa de aprendizado inicial igual a 0,001, e a função de perda categórica, adequada para
problemas de classificação multiclasse.

Para aprimorar o processo de treinamento e evitar sobreajuste, foram empregados mecanismos de controle por meio de callbacks. O ajuste dinâmico
da taxa de aprendizado foi realizado utilizando o método ReduceLROnPlateau, que reduz a taxa de aprendizado quando a função de perda de 
validação não apresenta melhora após um número definido de épocas. Além disso, foi adotado o critério de parada antecipada (Early Stopping), 
interrompendo o treinamento quando não há melhora significativa no desempenho de validação, restaurando automaticamente os melhores pesos 
obtidos. Também foi implementado o salvamento do melhor modelo durante o treinamento, com base no menor valor da perda de validação.

O treinamento foi realizado por até 100 épocas, com tamanho de lote igual a 64 amostras, utilizando os conjuntos de treinamento e teste 
definidos na etapa de pré-processamento. Ao final do treinamento, o modelo com melhor desempenho foi carregado e utilizado para a avaliação 
final, permitindo a obtenção das métricas de desempenho e da matriz de confusão para cada classe.

A definição dos hiperparâmetros do modelo, como o número de filtros, tamanhos de kernel, taxa de dropout e taxa de aprendizado, não foi realizada por 
meio de uma busca sistemática (como grid search ou random search). Esses valores foram escolhidos com base em práticas consolidadas na literatura para 
aplicações de redes convolucionais unidimensionais em sinais biomédicos, bem como em experimentações preliminares conduzidas neste trabalho.

A progressão do número de filtros (32, 64 e 128) foi adotada com o objetivo de permitir a extração hierárquica de características, partindo de padrões 
mais simples até representações mais complexas. Os tamanhos de kernel (5 e 3) foram definidos de forma a capturar padrões locais relevantes no sinal de 
ECG, como variações associadas ao complexo QRS e às ondas P e T, mantendo um equilíbrio entre capacidade de representação e custo computacional.

A taxa de dropout de 50\% foi empregada como uma estratégia de regularização, visando reduzir o risco de sobreajuste, especialmente considerando o uso 
de dados aumentados artificialmente. Já a taxa de aprendizado inicial de 0,001 foi definida com base na configuração padrão do otimizador Adam, sendo 
utilizada na literatura por apresentar boa estabilidade durante o treinamento.

Destaca-se que, embora esses valores tenham se mostrado adequados para o problema proposto, a ausência de um processo sistemático de otimização de 
hiperparâmetros constitui uma limitação deste trabalho, podendo impactar o desempenho final do modelo.

\section{Arquitetura e treinamento da CNN2D}

Com o objetivo de explorar informações espectro-temporais dos sinais de eletrocardiograma, foi desenvolvido um segundo modelo baseado em 
CNN2D. Diferentemente da abordagem unidimensional, nessa estratégia os sinais de ECG foram previamente transformados em representações 
bidimensionais por meio da análise tempo-frequência, permitindo que o modelo explorasse padrões espaciais associados às variações espectrais 
do sinal.

Inicialmente, cada segmento de ECG contendo 360 amostras foi convertido em uma imagem por meio da Transformada de Fourier de Curta Duração (STFT). Para 
essa transformação, foi utilizada frequência de amostragem de 360 Hz, tamanho de janela (nperseg) de 64 amostras e sobreposição (noverlap) de 32 
amostras, correspondente a 50\% de sobreposição entre janelas consecutivas. A implementação foi realizada por meio da função spectrogram da 
biblioteca SciPy, que utiliza, por padrão, janela do tipo Hamming.

A partir da STFT, foi obtido o espectrograma do sinal. Em seguida, foi aplicada uma transformação logarítmica na escala de potência, 
utilizando $10 \log_{10}(Sxx + 10^{-10})$, com o objetivo de realçar componentes de baixa intensidade e melhorar o contraste da imagem.

Em seguida, os valores foram normalizados no intervalo [0, 1] por meio de normalização min-max. Por fim, os espectrogramas foram redimensionados para o 
formato 32 × 32 utilizando interpolação bicúbica, sendo organizados como imagens com um único canal (32 × 32 × 1), utilizadas como entrada da CNN2D. A 
escolha dessa resolução representa um compromisso entre a preservação das características espectro-temporais relevantes e a eficiência computacional. 
Resoluções maiores, como 64 × 64 ou 128 × 128, implicariam aumento significativo no custo computacional e no número de parâmetros do modelo, sem 
garantia de ganhos proporcionais de desempenho, especialmente considerando o tamanho reduzido dos segmentos de entrada. Um exemplo pode ser 
visto na figura 3.

\begin{figure}[ht]
  \caption{Exemplo de uma transformação}
  \label{fig:exemplo2d}
  \centering
  \includegraphics[width=0.7\textwidth]{Exemplo_Classe_N_2D.png}
  \source
\end{figure}

O primeiro bloco convolucional é formado por uma camada Conv2D com 32 filtros, kernel de tamanho 3 x 3, função de ativação ReLU e preenchimento 
do tipo same, preservando as dimensões espaciais da imagem. Em seguida, foi aplicada a normalização em lote (Batch Normalization), com o 
objetivo de estabilizar o processo de treinamento. A redução dimensional foi realizada por meio de uma camada de MaxPooling2D com fator 2 x 2.

O segundo bloco convolucional segue estrutura semelhante, contendo uma camada Conv2D com 64 filtros e kernel 3 x 3, também utilizando ativação 
ReLU e padding same. Após a convolução, foram aplicadas novamente as etapas de normalização em lote e MaxPooling2D, permitindo a extração de 
padrões espaciais mais complexos associados às variações espectrais do ECG.

O terceiro bloco convolucional é composto por uma camada Conv2D com 128 filtros e kernel 3 x 3, seguida por Batch Normalization e MaxPooling2D. 
Esse bloco aprofunda a capacidade de representação do modelo, possibilitando a identificação de características discriminantes mais abstratas 
entre as classes de batimentos cardíacos.

Após os blocos convolucionais, a saída do modelo foi convertida em um vetor unidimensional por meio da camada Flatten. Em seguida, foi aplicada 
uma camada de Dropout com taxa de 50\%, com o objetivo de reduzir o sobreajuste. O classificador final é composto por uma camada densa com 
100 neurônios e função de ativação ReLU, seguida por uma camada de saída com cinco neurônios e função de ativação Softmax, correspondentes às 
cinco classes definidas pelo padrão AAMI.

\begin{table}[ht]
  \caption{Detalhamento da arquitetura proposta}
  \label{tab: table4}
  \centering

  \begin{tabular}{cccc}
    \hline
    \bfseries Camada & \bfseries Filtros/Neurônios & \bfseries Tamanho do Kernel & \bfseries Saída\\
    \hline
    Entrada	&	-	&	-	&	(32, 32, 1)\\
    Conv2D 1	&	32	&	3 x 3	&	(32, 32, 32)\\
    Batch Norm 1	&	-	&	-	&	(32, 32, 32)\\
    Max Pooling 1	&	-	&	2 x 2	&	(16, 16, 32)\\
    Conv2D 2	&	64	&	3 x 3	&	(16, 16, 64)\\
    Batch Norm 2	&	-	&	-	&	(16, 16, 64)\\
    Max Pooling 2	&	-	&	2 x 2	&	(8, 8, 64)\\
    Conv2D 3	&	128	&	3 x 3	&	(8, 8, 128)\\
    Batch Norm 3	&	-	&	-	&	(8, 8, 128)\\
    Max Pooling 3	&	-	&	2 x 2	&	(4, 4, 128)\\
    Flatten	&	-	&	-	&	(2048)\\
    Dropout	&	-	&	-	&	(2048)\\
    Densa (FC)	&	100	&	-	&	(100)\\
    Saída (Softmax)	&	5	&	-	&	(5)\\

    \hline
  \end{tabular}

  \source
\end{table}

O treinamento do modelo foi realizado utilizando o otimizador Adam, com taxa de aprendizado inicial igual a 0,001, e a função de perda 
categorical cross-entropy, adequada para problemas de classificação multiclasse.

Para melhorar a convergência e evitar sobreajuste, foram empregados mecanismos de controle baseados em callbacks. O método ReduceLROnPlateau 
foi utilizado para reduzir automaticamente a taxa de aprendizado quando a perda de validação não apresentava melhora após um número 
definido de épocas. Além disso, foi adotada a técnica de parada antecipada (Early Stopping), interrompendo o treinamento quando não havia 
ganho significativo no desempenho de validação, restaurando os melhores pesos obtidos.

O treinamento foi realizado por até 200 épocas, com tamanho de lote igual a 64 amostras, utilizando os conjuntos de treinamento e validação
previamente definidos. Ao final do processo, o modelo treinado foi utilizado para gerar as previsões no conjunto de teste, possibilitando a 
construção da matriz de confusão e o cálculo das métricas de desempenho para cada classe.

A definição dos hiperparâmetros do modelo, como número de filtros, tamanho dos kernels, taxa de dropout e taxa de aprendizado, não foi realizada por 
meio de uma busca sistemática. Esses valores foram definidos com base em práticas consolidadas na literatura para redes convolucionais aplicadas a 
imagens, bem como em experimentações preliminares conduzidas neste trabalho.

A progressão do número de filtros (32, 64 e 128) foi adotada com o objetivo de permitir a extração hierárquica de características espaciais, partindo 
de padrões mais simples até representações mais abstratas. O tamanho de kernel 3 × 3 foi escolhido por ser amplamente utilizado em arquiteturas 
convolucionais, apresentando bom desempenho na captura de padrões locais em imagens.

A taxa de dropout de 50\% foi empregada como estratégia de regularização, visando reduzir o risco de sobreajuste. A taxa de aprendizado inicial de 
0,001 foi definida com base na configuração padrão do otimizador Adam, sendo reconhecida por sua estabilidade em diferentes tarefas de aprendizado 
profundo.

Destaca-se que a ausência de um processo sistemático de otimização de hiperparâmetros constitui uma limitação deste trabalho, podendo impactar o 
desempenho final do modelo.